\documentclass[11pt]{article}
\usepackage[T1]{fontenc}
\usepackage[margin=1in]{geometry}
\usepackage{amsmath}
\usepackage[breakable,skins]{tcolorbox}

\tcbset{
  colback=white,
  colframe=blue!60!black,
  boxrule=0.8pt,
  fonttitle=\bfseries,
  left=4pt, right=4pt, top=4pt, bottom=4pt,
}

\newtcolorbox{keypoint}{
  colback=yellow!10,
  colframe=orange!70!black,
  title=Key point,
}

\title{Payments Service: Incident Runbook Excerpts}
\author{Platform Engineering}
\date{October 2026}

\begin{document}
\maketitle

\section{Overview}
This document collects the boxed callouts used in the payments service
runbook. Every box follows the house style set with \verb|\tcbset|:
a white body, a dark blue frame, and a bold title. Admonitions that
apply to the whole page use the custom \verb|keypoint| environment,
which renders with an amber frame so it stands out from ordinary
notes.

\begin{tcolorbox}[title=Scope]
These excerpts cover the deploy freeze calendar, the rollback
checklist, and the on-call escalation path. They do not cover
database failover, which has its own runbook.
\end{tcolorbox}

\begin{keypoint}
During the holiday freeze---the last two weeks of December---no
production deploys are allowed without written approval from two
service owners.
\end{keypoint}

\section{Escalation}
The escalation path below uses the \verb|skins| library's enhanced
skin with a soft drop shadow, which the team reserves for boxes that
must catch the eye during an incident.

\begin{tcolorbox}[enhanced, title=Escalation path, drop fuzzy shadow,
  colback=red!5, colframe=red!60!black]
Page the primary on-call first. If there is no acknowledgement within
ten minutes, page the secondary. If the incident is customer-facing
and unresolved after thirty minutes, declare a major incident and
summon the incident commander.
\end{tcolorbox}

\section{Rollback Checklist}
The rollback checklist is the longest box in the runbook: it usually
runs past the page break, so it is declared breakable and allowed to
split across pages.

\begin{tcolorbox}[breakable, title=Rollback checklist (v14)]
Confirm that the deploy pipeline is paused before touching anything;
a second deploy landing mid-rollback has caused two of the last five
rollback incidents. Announce the rollback in the incidents channel
with the release number, the reason, and the expected duration, so
that support can set customer expectations early.

Drain traffic from the new revision one availability zone at a time,
waiting for the health checks to go green on the old revision before
moving to the next zone. The load balancer target groups are named
\texttt{payments-blue} and \texttt{payments-green}; the runbook
assumes blue is live, but verify against the deploy ticket rather
than trusting the colour, because the colours swap on every release.

Watch the error budget dashboard while each zone drains. A brief
spike in 5xx responses is normal as connections migrate, but if the
burn rate stays above the fast-burn threshold for more than five
minutes, stop the drain and escalate: the old revision may itself be
unhealthy, and completing the rollback would only move the incident.

Once all zones serve the old revision, run the smoke suite from the
operations host, not from a laptop: the suite needs the production
VPN routes, and the VPN client on laptops drops idle connections
after fifteen minutes. The suite takes about six minutes; all
forty-two checks must pass before the rollback is declared complete.

File the rollback report within one business day. The report needs
the timeline, the graphs, and a link to the deploy diff, plus a
judgement on whether the release can be re-attempted or must be
reverted in source control. Attach the report to the original deploy
ticket so the history stays in one place.

Finally, schedule the retrospective within a week while memories are
fresh. The template is in the shared drive under \emph{Engineering /
Incident templates}. Invite the deploy author, the reviewer, and the
on-call who performed the rollback; incidents without a
retrospective within two weeks are flagged in the engineering review.
\end{tcolorbox}

\begin{keypoint}
A rollback is complete only when the smoke suite passes on the old
revision and the rollback report is filed. Draining traffic is
necessary but not sufficient.
\end{keypoint}
\end{document}
